\documentclass[11pt,a4paper]{article}
\usepackage[margin=25mm]{geometry}
\usepackage{amsmath,amssymb,booktabs,siunitx,hyperref,listings,microtype}
\hypersetup{colorlinks=true,linkcolor=blue,urlcolor=blue,citecolor=blue,pdftitle={Report 2 - Rigid Pipe Model for OpenHPL/OpenHPLjl},pdfauthor={Madhusudhan Pandey}}
\lstdefinelanguage{Julia}{keywords={using,begin,end,function,return,if,else,for,while,module,export,include},sensitive=true,comment=[l]\#,morestring=[b]"}
\lstset{language=Julia,basicstyle=\ttfamily\small,breaklines=true,frame=single,columns=fullflexible,showstringspaces=false}
\title{Report 2 - Rigid Pipe Model for OpenHPL/OpenHPLjl}
\author{Madhusudhan Pandey}
\date{19 September 2026}
\begin{document}
\maketitle
\begin{abstract}
This report introduces the rigid-pipe component as the second hydraulic model in OpenHPLjl. The model adds momentum conservation and water inertia to the reservoir storage model developed previously. The same documentation sequence is retained for consistency: context, generalized concept, literature alternatives, mathematics, OpenHPL interpretation, ModelingToolkit implementation, assembly, validation, expected outputs, and engineering interpretation.
\end{abstract}

\section{Updated Context}
The reservoir establishes hydraulic potential and storage. The next component must describe how water accelerates through a conduit when a head difference is applied. For a rigid, incompressible water column, the dominant dynamic is lumped water inertia:
\[
\boxed{\text{Reservoir}\rightarrow\text{RigidPipe}\rightarrow\text{SurgeTank}\rightarrow\text{Penstock}\rightarrow\text{Turbine}}.
\]
This model is appropriate when elastic water-hammer effects are not the primary phenomenon of interest. It is especially useful for governor, AGC, and FCR studies where a compact nonlinear waterway model is preferred.

\section{Generalized Concept Model}
A rigid pipe has two hydraulic ports and one principal dynamic quantity, the volumetric flow $Q$:
\[
\boxed{(H_{\rm in},Q_{\rm in})\rightarrow\text{water inertia + losses}\rightarrow(H_{\rm out},Q_{\rm out})}.
\]
With the OpenHPLjl connector convention, flow is positive into each component. Therefore
\[
Q_{\rm in}=Q,\qquad Q_{\rm out}=-Q,
\]
and mass conservation gives
\[
Q_{\rm in}+Q_{\rm out}=0.
\]

\section{Other Models Available in the Literature}
Common waterway representations can be ordered by fidelity:
\[
\boxed{\text{quasi-steady loss}\rightarrow\text{rigid water column}\rightarrow\text{elastic lumped model}\rightarrow\text{method of characteristics / PDE}}.
\]
The quasi-steady model neglects inertia. The rigid-water-column model retains one flow state. Elastic models add pipe-wall and fluid compressibility. Distributed models resolve pressure waves spatially and are preferred for water hammer.

\section{Mathematics Involved in the OpenHPL Concept/Model}
For a constant-area conduit of length $L$ and cross-section $A$, a lumped momentum balance can be written in head form as
\[
\frac{L}{gA}\frac{dQ}{dt}=H_{\rm in}-H_{\rm out}-H_f(Q).
\]
Define the hydraulic inertance
\[
T_w=\frac{L}{gA}.
\]
Using a quadratic loss law
\[
H_f(Q)=R Q|Q|,
\]
the model becomes
\[
\boxed{T_w\dot Q=H_{\rm in}-H_{\rm out}-RQ|Q|}.
\]
Equivalently,
\[
\boxed{\dot Q=\frac{gA}{L}\left(H_{\rm in}-H_{\rm out}-RQ|Q|\right)}.
\]
The state is $x=Q$. Hydraulic heads at the two ports and the port flows participate in algebraic connection equations.

\section{OpenHPL-Specific Modeling Interpretation}
OpenHPL waterway models are based on conservation of mass and momentum. Its conduit formulation includes pressure-gradient, gravity, inertia, and friction terms. The rigid-pipe model used here is a reduced head-based representation of the same physical balance, chosen as the first reusable OpenHPLjl waterway element. More detailed elastic or distributed formulations can later preserve the same hydraulic connector.

\section{Implementation in Julia ModelingToolkit / SciML}
The component reuses the existing \texttt{HydraulicPort}. ModelingToolkit's acausal connection semantics equate head variables and sum \texttt{Flow}-marked quantities to zero.

\begin{lstlisting}
@component function RigidPipe(; name,
    L = 100.0,
    A = 1.0,
    g = 9.81,
    R = 0.0,
    Q0 = 0.0)

    @named inlet = HydraulicPort()
    @named outlet = HydraulicPort()

    @parameters L=L A=A g=g R=R
    @variables Q(t)=Q0

    eqs = [
        inlet.Q ~ Q
        outlet.Q ~ -Q
        D(Q) ~ (g*A/L) *
            (inlet.H - outlet.H - R*Q*abs(Q))
    ]

    System(eqs, t, [Q], [L,A,g,R];
        systems=[inlet,outlet], name=name)
end
\end{lstlisting}

\section{Model Assembly and Connections}
The component is connected without assigning causal input/output directions:
\begin{lstlisting}
@named reservoir = InfiniteReservoir(H0 = 100.0)
@named pipe = RigidPipe(L = 100.0, A = 1.0)

connections = [
    connect(reservoir.port, pipe.inlet)
]
\end{lstlisting}
A downstream boundary, surge tank, valve, or turbine supplies the outlet relation.

\section{Numerical Experiment / Validation}
For a frictionless pipe,
\[
L=100~\si{\meter},\quad A=1~\si{\meter\squared},\quad g=9.81~\si{\meter\per\second\squared},
\]
with
\[
H_{\rm in}=100~\si{\meter},\qquad H_{\rm out}=90~\si{\meter},\qquad Q(0)=0,
\]
we obtain
\[
\dot Q=\frac{9.81\cdot1}{100}(100-90)=0.981~\si{\meter\cubed\per\second\squared}.
\]
Hence after one second,
\[
\boxed{Q(1)=0.981~\si{\meter\cubed\per\second}}.
\]
This analytical result is the first unit-test target. A second test can introduce $R>0$ and verify convergence toward the steady solution
\[
Q_{\rm ss}=\sqrt{\frac{H_{\rm in}-H_{\rm out}}{R}}
\]
for positive flow.

\section{Expected Outputs}
The minimum useful outputs are
\[
Q(t),\qquad H_{\rm in}(t),\qquad H_{\rm out}(t),\qquad \Delta H(t),\qquad H_f(t).
\]
Recommended plots are flow response to a head step, head difference, friction loss, and the momentum-balance residual. For the frictionless validation case, $Q(t)$ should be a straight ramp with slope $0.981~\si{\meter\cubed\per\second\squared}$.

\section{Engineering Interpretation}
The rigid pipe introduces the second fundamental hydraulic dynamic after reservoir storage:
\[
\boxed{\text{Reservoir: mass/storage}\quad\rightarrow\quad\text{RigidPipe: momentum/inertia}}.
\]
It is deliberately simple, transparent, and reusable. The same interface can later be replaced by elastic penstock or distributed water-hammer models when higher-frequency pressure-wave dynamics are required.

\section*{Conclusion}
The first OpenHPLjl waterway state is the pipe flow $Q$. Its governing equation is
\[
\boxed{\frac{L}{gA}\dot Q=H_{\rm in}-H_{\rm out}-RQ|Q|}.
\]
The next component is the surge tank, which couples waterway momentum with local storage and introduces the first hydraulic oscillatory mode.

\begin{thebibliography}{9}
\bibitem{OpenHPL}
OpenHPL 3.0.1 documentation and resources.\newline
\url{https://build.openmodelica.org/Documentation/OpenHPL.html}

\bibitem{Sharefi}
S. Sharefi, hydropower conduit and surge-shaft momentum formulations, OpenHPL resources.\newline
\url{https://build.openmodelica.org/Documentation/OpenHPL%203.0.1/Resources/Documents/Sharefi2011.pdf}

\bibitem{MTK}
SciML, ModelingToolkit.jl, Acausal Component-Based Modeling.\newline
\url{https://docs.sciml.ai/ModelingToolkit/stable/tutorials/acausal_components/}
\end{thebibliography}
\end{document}
